% !TEX root = ../main.tex
\begin{abstract}
Long-term timber harvest schedules produced by linear programming (LP) are routinely used in decision support. They are presented to decision-makers as credible statements about the future, and their shadow prices are used to value policy trade-offs. An open-loop LP solution is a plan made as if the planner could precommit to every future decision. However, real planning is sequential and adaptive: each period, the planner observes the realized forest state and re-solves. In a 1991 PhD thesis, P. J. Daugherty showed that LP forest plans solved open-loop admit \emph{dynamically inconsistent} plans. These plans fail Bellman's principle of optimality because a future planner re-solving from the realized state would not follow the plan's tail. We formalize dynamic inconsistency for LP forest plans with an operational divergence metric and provide the first open, reproducible benchmark of the phenomenon, reproducing and refining Daugherty's Umpqua National Forest FORPLAN case study in the open-source \texttt{ws3} model (Paradis 2026) from public planning documents. Across a grid of tested scenarios (18 landbases $\times$ 4 discount rates $\times$ 6 harvest-flow policies), dynamic inconsistency is pervasive under harvest-flow constraints (73\% of scenarios) yet largely absent without them at positive discount rates. The tested scenarios confirm that the formulation of the harvest volume flow constraint is the operative between-period link which drives the inconsistency. Inconsistency occurs at every discount rate, its magnitude is rate-dependent, and the non-declining-yield policy is the most inconsistent. On a landbase dominated by financially mature stands, the open-loop non-declining-yield plan promises a harvest flow that sequential replanning revises downward over most of the horizon, realizing today's harvests at the cost of future productive potentials. Four extension experiments test whether the phenomenon can be designed away: declining discount-rate schedules make it strictly more pervasive and introduce a preference-level (Strotz) inconsistency channel of their own; denominating the flow constraint in net revenue rather than volume doubles its magnitude while proving the negatively valued strata are a symptom rather than the cause; a rolling-mean flow floor shows that the replanning institution, not the constraint's shape, is the operative margin; and a max-harvest cap calibrated by iterative re-solve to a realized even-flow criterion eliminates the inconsistency entirely---at a credible harvest level measurably below the announced one. The presented model, data, and experiment are openly available and contribute to increased awareness of the risk of dynamically inconsistent LP solutions under even-flow constraints at the intersection of forest management planning, forest economics research, and decision support.
\end{abstract}
